\begin{figure}[!ht]
  \centering
  \begin{tikzpicture}[alt={A native Memory row contains the first sixteen inflated bytes. An outline marks six Up residual bytes, all 05. Dashed outlines identify the preceding None and Up selectors.}]
    \node[anchor=north west,inner sep=0,font=\small\bfseries] (title) at (0,0) {The Up row still contains residuals};
    \node[anchor=north west,inner sep=0] (a) at ([yshift=-2mm]title.south west) {\MARSRevisionCrop{revision/captures/png/png-inflated.png}{0.99\linewidth}{0.03125}{0.33888889}{0.58203125}{0.60888889}};
    \begin{scope}[shift={(a.south west)},x={($(a.south east)-(a.south west)$)},y={($(a.north west)-(a.south west)$)}]
      \draw[black,line width=.7pt] (.5825,.01) rectangle (.869375,.467);
      \draw[black,dashed,line width=.7pt] (.20,.01) rectangle (.2478125,.467);
      \draw[black,dashed,line width=.7pt] (.5346875,.01) rectangle (.5825,.467);
    \end{scope}
  \end{tikzpicture}
  \caption{Before reconstruction, the six outlined bytes are Up residuals of five. The dashed selectors are zero for the first row and two for the second. The final two pictured bytes begin the next row; they are outside the outlined Up samples.}
  \label{fig:revision-png-native-inflated}
\end{figure}
